\subsection{Best-of-\texorpdfstring{$K$}{K} Phase-Rescue Diagnostic}
\label{sec:phase_rescue}

The best-of-$K$ coverage argument in \refsec{sec:method} provides intuition
about candidate coverage under independent sampling; it does not itself
establish recovery of phase-compatible candidates. We therefore use an
offline paired diagnostic on held-out expert demonstrations to isolate
this mechanism.

For each environment, the first eight demonstrations form the reference
bank and the next 32 (zero-based demonstration indices 8--39) provide
query windows. Beyond the eight demonstrations used for training, these
32 additional demonstrations are used only for this offline diagnostic;
they are not supplied to policy optimization. We sample 512 query windows
without replacement from all valid windows (stride one), and construct
four candidate pools per query, giving 2048 paired comparisons. Each
window contains 32 state transitions, represented by concatenated
successive observations. We use unstandardized transition features for
OT with cosine ground cost, Sinkhorn regularization $0.1$, and 60
iterations. The $K=1$ condition receives a reference window sampled
uniformly from those whose circular phase error is within $0.04$ cycles
of the requested mismatch. The $K=8$ condition receives that same window
plus seven reference windows sampled uniformly with replacement.
The requested mismatch here is $0.5$ cycles; discretization and the
sampling tolerance mean the realized mismatch need not equal $0.5$.

\paragraph{Period and shared phase reference.}
Phase estimation uses a separate standardized copy of the reference
observations. Means and population standard deviations are computed
across all valid timesteps of the eight reference demonstrations;
dimensions with standard deviation at most $10^{-4}$ are excluded.
For each integer lag $\ell$, we compute the mean squared difference
between standardized observations at $t$ and $t+\ell$, averaging first
over valid timesteps and dimensions within each reference demonstration
and then equally over the eight demonstrations. We discard the first
100 timesteps for this calculation. The inclusive lag ranges are
$[10,90]$ for Hopper, $[8,80]$ for Ant, and $[8,120]$ for both humanoids.
We select the strict interior local minimum with the lowest score;
if none exists, we select the lowest-scoring lag over the full range.
This yields the period $P$ in simulation steps.

For each environment, we take observations at zero-based indices
$100,\ldots,100+P-1$ of the first reference demonstration as the cycle
template. This is \emph{one shared template per environment}, not a
separate template for each demonstration. Each reference window is
assigned a discrete phase $p(c)\in\{0,\ldots,P-1\}$ by comparing its
first eight standardized observations with the eight observations
starting at each circular offset of the template (wrapping modulo $P$)
and choosing the offset with minimum squared Euclidean distance.
The query target phase $p^\star(q)$ is the phase of its nearest reference
window under time-aligned squared Euclidean distance over the first
eight unstandardized transition vectors, using all observation coordinates
in both states of each transition and searching the entire
reference bank independently of the sampled pool and OT scores.
Ties are resolved by the first minimum in index order.

\paragraph{Metrics and uncertainty.}
The circular phase error is
\[
 e(q,c)=\frac{\min\{|p(c)-p^\star(q)|,
 P-|p(c)-p^\star(q)|\}}{P}\in[0,0.5].
\]
A candidate is phase-compatible if $e(q,c)\leq0.10$ cycles.
Coverage is the percentage of augmented pools containing at least one
phase-compatible candidate.
Let $c_1$ be the shared candidate and $\hat c_8$ the candidate with
minimum OT cost in the augmented pool. A rescue occurs when
$e(q,c_1)>0.10$ and $e(q,\hat c_8)\leq0.10$; phase improvement means
$e(q,\hat c_8)<e(q,c_1)$. Coverage, rescue, and improvement percentages
use all 2048 pairs as their denominator. Mean phase errors describe
the minimum-cost candidate, not the soft-weighted mixture.
We report pointwise 95\% percentile confidence intervals from 5000
bootstrap draws of query trajectories with replacement. The cluster
population consists of the query trajectories with at least one sampled
window. In these runs, all 32 held-out trajectories (demonstration indices
8--39) contributed windows; each bootstrap draw therefore samples 32
trajectory IDs with replacement. A trajectory with no sampled windows
would be excluded from the cluster population. All windows
and pool repetitions from a sampled trajectory remain together;
statistics are recomputed over the resulting pairs. The reference
bank and phase template remain fixed, so these intervals quantify
query-trajectory sampling uncertainty conditional on those references.
The sampling seed is 0 and the bootstrap seed is 12345.
For Hopper, Ant, and Humanoid, the discrete phases and the mismatch
interval $[0.46,0.50]$ admit only one circular mismatch magnitude:
$18/37$, $11/23$, and $10/20$, respectively. Their $K=1$ errors are
therefore fixed by construction, giving zero-width bootstrap intervals
(up to floating-point precision). SNU Humanoid admits both $14/30$ and
$15/30$, so its baseline error can vary.

\begin{table}[H]
\centering\small
\setlength{\tabcolsep}{0.35em}
\begin{tabular}{lccc}
\toprule
Environment & Coverage (\%) & Candidate rescued (\%) & Phase improved (\%) \\
\midrule
Hopper & 78.2 [76.5, 79.9] & 46.3 [43.6, 48.9] & 96.8 [95.9, 97.7] \\
Ant & 82.0 [80.2, 83.8] & 63.4 [60.9, 65.9] & 98.5 [98.0, 99.1] \\
Humanoid & 85.2 [83.8, 86.7] & 64.0 [62.1, 65.9] & 97.1 [96.5, 97.7] \\
SNU Humanoid & 85.1 [83.8, 86.5] & 72.0 [70.1, 73.9] & 99.3 [99.0, 99.6] \\
\bottomrule
\end{tabular}
\caption{Paired phase-rescue diagnostic at a target half-cycle mismatch.
Entries are point estimates [95\% trajectory-cluster bootstrap CI].
All percentages use all 2048 query--pool pairs per environment.}
\label{tab:phase_rescue}
\end{table}
\begin{table}[H]
\centering\small
\begin{tabular}{lrcc}
\toprule
Environment & Period $P$ (steps) & Phase error $K=1$ & Phase error $K=8$ \\
\midrule
Hopper & 37 & 0.486 [0.486, 0.486] & 0.147 [0.139, 0.156] \\
Ant & 23 & 0.478 [0.478, 0.478] & 0.102 [0.098, 0.107] \\
Humanoid & 20 & 0.500 [0.500, 0.500] & 0.121 [0.116, 0.125] \\
SNU Humanoid & 30 & 0.477 [0.477, 0.478] & 0.095 [0.092, 0.099] \\
\bottomrule
\end{tabular}
\caption{Mean circular phase errors in gait cycles [95\% trajectory-cluster
bootstrap CI]. Phase labels use a single shared reference-cycle template
per environment.}
\label{tab:phase_rescue_errors}
\end{table}


\paragraph{Effect of candidate-pool size.}
We extend the half-cycle diagnostic to $K\in\{4,8,16\}$ using nested
candidate pools and the same 2048 query--pool pairs per environment.
The $K=4$ pool consists of the first four candidates of each verified
$K=8$ pool, including the shared mismatched candidate. The $K=16$ pool
appends eight independent, uniformly sampled reference windows with
replacement (additional-candidate seed 67890). We retain the original
query windows, phase labels, and $K=8$ pools, and verify their OT costs
by recomputation. Thus the $K=8$ results are unchanged.
Coverage and rescue retain the same definitions with the selected
candidate taken from the corresponding $K$-candidate pool.

\begin{table}[H]
\centering\small
\begin{tabular}{lrccc}
\toprule
Environment & $K$ & Coverage (\%) & Rescue (\%) & Mean phase error \\
\midrule
hopper & 4 & 47.3 [44.9, 49.7] & 32.6 [30.2, 35.0] & 0.200 [0.190, 0.211] \\
hopper & 8 & 78.2 [76.5, 79.9] & 46.3 [43.6, 48.9] & 0.147 [0.139, 0.156] \\
hopper & 16 & 95.8 [94.9, 96.6] & 56.6 [54.1, 58.8] & 0.118 [0.112, 0.124] \\
ant & 4 & 51.8 [49.8, 53.9] & 40.1 [38.0, 42.3] & 0.168 [0.162, 0.174] \\
ant & 8 & 82.0 [80.2, 83.8] & 63.4 [60.9, 65.9] & 0.102 [0.098, 0.107] \\
ant & 16 & 97.7 [97.2, 98.3] & 81.4 [79.6, 83.3] & 0.073 [0.069, 0.078] \\
humanoid & 4 & 56.1 [54.4, 57.9] & 43.0 [41.3, 44.8] & 0.190 [0.184, 0.196] \\
humanoid & 8 & 85.2 [83.8, 86.7] & 64.0 [62.1, 65.9] & 0.121 [0.116, 0.125] \\
humanoid & 16 & 98.0 [97.4, 98.7] & 77.7 [76.0, 79.3] & 0.089 [0.085, 0.093] \\
snu\_humanoid & 4 & 55.8 [53.7, 58.0] & 48.0 [45.9, 50.4] & 0.168 [0.163, 0.174] \\
snu\_humanoid & 8 & 85.1 [83.8, 86.5] & 72.0 [70.1, 73.9] & 0.095 [0.092, 0.099] \\
snu\_humanoid & 16 & 97.9 [97.4, 98.5] & 87.2 [85.6, 88.8] & 0.058 [0.055, 0.062] \\
\bottomrule
\end{tabular}
\caption{Nested-pool phase-rescue comparison at a target half-cycle mismatch. Entries are point estimates [95\% trajectory-cluster bootstrap CI], using 512 held-out query windows and four pools per window. $K=4$ uses the first four candidates of the verified $K=8$ pool; $K=16$ appends eight uniformly sampled candidates. All conditions share the same mismatched baseline.}
\label{tab:phase_rescue_k_sweep}
\end{table}


Increasing $K$ improves coverage and rescue rates and reduces mean
selected phase error in all four environments. From $K=8$ to $K=16$,
rescue increases by 10.3, 18.0, 13.7, and 15.1 percentage points for
Hopper, Ant, Humanoid, and SNU Humanoid, respectively. Paired 95\%
trajectory-cluster bootstrap intervals for these increases are
$[8.5,12.3]$, $[16.4,19.7]$, $[12.0,15.4]$, and $[13.2,17.1]$ percentage
points. These intervals are computed on within-pair differences,
using the same trajectory-cluster resampling procedure, and are
pointwise without multiplicity correction. The phase-error reduction
from $K=8$ to $K=16$ is smaller than from $K=4$ to $K=8$ in every
environment. These results concern minimum-cost selection on held-out
expert windows, not policy-training performance.


This diagnostic measures recovery of phase-compatible candidates in
held-out expert windows. It does not directly measure policy return,
gradient alignment, or the effect of the soft-weighted training objective,
and does not guarantee robustness when compatible candidates are absent
or under arbitrary distribution shift.
